\exercise{Traffic: Spectral centrality}{3}
Consider once again the road map from the previous exercise, but now compute the spectral centrality. (Mind the parallel links.)

\solution
We represent the parallel links by setting the corresponding elements of the adjacency matrix to 2, like we did in Chap.~3. The resulting matrix is 
\eq{
{\bf A} =  \left(\begin{array}{ c c c c c c c c c} 
0 & 1 & 1 & 0 & 0 & 0 & 0 & 0 & 1 \\
1 & 0 & 1 & 1 & 0 & 0 & 1 & 0 & 0 \\
1 & 1 & 0 & 0 & 2 & 0 & 0 & 0 & 0 \\
0 & 1 & 0 & 0 & 0 & 0 & 1 & 1 & 0 \\
0 & 0 & 2 & 0 & 0 & 1 & 0 & 1 & 0 \\
0 & 0 & 0 & 0 & 1 & 0 & 2 & 1 & 0 \\
0 & 1 & 0 & 1 & 0 & 2 & 0 & 0 & 0 \\
0 & 0 & 0 & 1 & 1 & 1 & 0 & 0 & 0 \\
1 & 0 & 0 & 0 & 0 & 0 & 0 & 0 & 0
\end{array}\right) 
}
where I used the ordering 
\begin{center}
A, B, C, D, E, F, G, H, X  
\end{center}
I had to multiply my starting vector about 10 times until it stabilized. I then normalized the vector such that the first entry is (the one corresponding to node A) is 1. The relative importance of the nodes is then:
\begin{center}
A: 1, B: 1.6, C: 1.7, D: 1.3, E: 1.8, F: 1.9, G: 1.8, H: 1.4, X: 0.2
\end{center}
Note that the least important node is again $X$, while the most important is now F. 
